\section{Trabajo relacionado}
\label{sec:relacionado}

\paragraph{Conjuntos de datos de gestos.}
HaGRID \citep{kapitanov2024hagrid} es uno de los mayores conjuntos de imágenes de gestos de la
mano: su primera versión contiene 552\,992 imágenes FullHD de 18 gestos pensados para controlar
dispositivos, tomadas por miles de personas en entornos domésticos y con condiciones de
iluminación muy variadas. Cada imagen incluye las cajas de las manos, la etiqueta de cada mano
(las manos que no hacen el gesto se etiquetan como \cls{no\_gesture}) y un identificador
anónimo de la persona (\code{user\_id}), lo que permite particionar por persona. Usamos un
subconjunto público de 30\,000 imágenes reescaladas a 384\,px \citep{mills_hagridsample}.

\paragraph{Detector seguido de clasificador.}
En una imagen completa la mano ocupa muy pocos píxeles, así que los sistemas prácticos suelen
dividir el problema en dos etapas: un detector o rastreador de manos, como MediaPipe Hands
\citep{zhang2020mediapipe}, localiza la mano, y un clasificador decide el gesto a partir del
recorte \citep{pisharady2015review}. Este trabajo se centra en la segunda etapa: durante el
entrenamiento y la evaluación usamos las cajas anotadas, y el detector solo interviene en la
demo (\cref{sec:demo}).

\paragraph{Arquitecturas convolucionales.}
Nuestra red sigue el esquema de bloques de VGG \citep{simonyan2015vgg}: dos convoluciones
$3\times3$ seguidas de \emph{max pooling}. Le añade normalización por lotes (BN)
\citep{ioffe2015bn}, que estabiliza la distribución de las activaciones y permite tasas de
aprendizaje mayores, y \emph{dropout} \citep{srivastava2014dropout}, en su variante espacial
(\code{Dropout2d}) en los bloques convolucionales \citep{tompson2015spatialdropout}. En lugar de
capas densas grandes al final usa \emph{global average pooling} \citep{lin2014nin}, que reduce
mucho el número de parámetros. Como referencia de un modelo más grande usamos ResNet18
\citep{he2016resnet} preentrenado en ImageNet \citep{deng2009imagenet}; ajustar una red
preentrenada aprovecha rasgos genéricos de las primeras capas que se transfieren bien a otras
tareas \citep{yosinski2014transferable}.

\paragraph{Explicabilidad y evaluación.}
Grad-CAM \citep{selvaraju2017gradcam} genera mapas de calor que indican qué regiones de la
imagen empujan la predicción hacia una clase, a partir de los gradientes en la última capa
convolucional. Lo usamos para comprobar que la red mira la mano y no el fondo. Para comparar
clasificadores evaluados sobre el mismo conjunto de test usamos la prueba de McNemar
\citep{mcnemar1947}, recomendada para este escenario por \citet{dietterich1998tests}, e
intervalos de confianza por \emph{bootstrap} \citep{efron1993bootstrap}.
